% Seminar 3 - Pagina 8. Fragment pentru
% % Seminar 3 - Pagina 8. Fragment pentru
% % Seminar 3 - Pagina 8. Fragment pentru
% % Seminar 3 - Pagina 8. Fragment pentru
% \input{seminars/pages/seminar3-pagina8-ro}; fără preambul sau pagină nouă.
% NOTE EDITORIALE:
% - Manuscrisul folosește Hom(A,B) pentru mulțimea tuturor aplicațiilor A -> B.
% - O clasă de echivalență este numită în notițe „germenele” lui f pe A'.
%   În terminologia uzuală, germenii presupun acordul pe o vecinătate;
%   relația de aici cere numai egalitatea restricțiilor la A'.
% - Formula exactă de după „germenele lui f” este neclară; pare să
%   trimită la A', posibil prin expresia „în punctul lui A'”.
% - Diagrama cu restricția este reprodusă folosind tikz-cd; p este
%   proiecția canonică, iar h aplicația de restricție.
% - Construcția de la final alege a în B, presupunând implicit B nevidă
%   dacă A\A' este nevidă; această ipoteză nu este enunțată în manuscris.
% - Semnele barate lângă diagrame nu sunt reconstituite.
% - Manuscrisul se oprește la enunțarea existenței unei prelungiri;
%   nu am adăugat o încheiere nouă a demonstrației.

\subsection*{Seminar 3 - Pagina 8}
\begingroup
\setlength{\abovedisplayskip}{5pt}
\setlength{\belowdisplayskip}{5pt}
\setlength{\abovedisplayshortskip}{3pt}
\setlength{\belowdisplayshortskip}{3pt}

Fie $A,B$ mulțimi și $A'\subseteq A$.
Pentru $f,g\in\operatorname{Hom}(A,B)$, definim
\[
 f\sim g\overset{\mathrm{def}}\iff f|_{A'}=g|_{A'}.
\]
Să se arate că $\sim$ este relație de echivalență și că
\[
 \operatorname{Hom}(A,B)/{\sim}
 \ \overset{\mathrm{bij}}{\cong}\ \operatorname{Hom}(A',B).
\]
Un element din mulțimea factor se numește în notițe
\emph{germenele} lui $f$ pe $A'$.

\medskip
\noindent\underline{Demonstrație.}
Considerăm proiecția canonică $p$ și aplicația de restricție $h$:
\[
\begin{tikzcd}[column sep=huge,row sep=large]
 \operatorname{Hom}(A,B)
    \arrow[r,"p"]\arrow[d,"h"']
  & \operatorname{Hom}(A,B)/{\sim}
       \arrow[dl,dashed,"\overline h"] \\
 \operatorname{Hom}(A',B) &
\end{tikzcd}
\qquad h(f)=f|_{A'}.
\]
Relația este cea indusă de $h$:
\[
 \sim\;=R_h,
 \qquad f\sim g\iff h(f)=h(g)
 \iff f|_{A'}=g|_{A'}.
\]
Prin urmare, aplicația indusă $\overline h$ este injectivă.

\medskip
\noindent Pentru orice aplicație $t:A'\longrightarrow B$ căutăm
$f:A\longrightarrow B$ astfel încât $f|_{A'}=t$.
Alegem $a\in B$ și definim
\[
 f(x)=
 \begin{cases}
   t(x),&x\in A',\\
   a,&x\in A\setminus A'.
 \end{cases}
\]
Aceasta este prelungirea indicată de diagramă:
\[
\begin{tikzcd}[column sep=large,row sep=large]
 A'\arrow[r,hook,"i"]\arrow[d,"t"']
    & A\arrow[dl,dashed,"f"]\\
 B &
\end{tikzcd}
\qquad f\circ i=t.
\]
Astfel, cum apare la sfârșitul paginii,
\[
 \forall\,g:A'\longrightarrow B,\quad
 \exists\,f:A\longrightarrow B
 \quad\text{a.î.}\quad f|_{A'}=g.
\]

\endgroup
; fără preambul sau pagină nouă.
% NOTE EDITORIALE:
% - Manuscrisul folosește Hom(A,B) pentru mulțimea tuturor aplicațiilor A -> B.
% - O clasă de echivalență este numită în notițe „germenele” lui f pe A'.
%   În terminologia uzuală, germenii presupun acordul pe o vecinătate;
%   relația de aici cere numai egalitatea restricțiilor la A'.
% - Formula exactă de după „germenele lui f” este neclară; pare să
%   trimită la A', posibil prin expresia „în punctul lui A'”.
% - Diagrama cu restricția este reprodusă folosind tikz-cd; p este
%   proiecția canonică, iar h aplicația de restricție.
% - Construcția de la final alege a în B, presupunând implicit B nevidă
%   dacă A\A' este nevidă; această ipoteză nu este enunțată în manuscris.
% - Semnele barate lângă diagrame nu sunt reconstituite.
% - Manuscrisul se oprește la enunțarea existenței unei prelungiri;
%   nu am adăugat o încheiere nouă a demonstrației.

\subsection*{Seminar 3 - Pagina 8}
\begingroup
\setlength{\abovedisplayskip}{5pt}
\setlength{\belowdisplayskip}{5pt}
\setlength{\abovedisplayshortskip}{3pt}
\setlength{\belowdisplayshortskip}{3pt}

Fie $A,B$ mulțimi și $A'\subseteq A$.
Pentru $f,g\in\operatorname{Hom}(A,B)$, definim
\[
 f\sim g\overset{\mathrm{def}}\iff f|_{A'}=g|_{A'}.
\]
Să se arate că $\sim$ este relație de echivalență și că
\[
 \operatorname{Hom}(A,B)/{\sim}
 \ \overset{\mathrm{bij}}{\cong}\ \operatorname{Hom}(A',B).
\]
Un element din mulțimea factor se numește în notițe
\emph{germenele} lui $f$ pe $A'$.

\medskip
\noindent\underline{Demonstrație.}
Considerăm proiecția canonică $p$ și aplicația de restricție $h$:
\[
\begin{tikzcd}[column sep=huge,row sep=large]
 \operatorname{Hom}(A,B)
    \arrow[r,"p"]\arrow[d,"h"']
  & \operatorname{Hom}(A,B)/{\sim}
       \arrow[dl,dashed,"\overline h"] \\
 \operatorname{Hom}(A',B) &
\end{tikzcd}
\qquad h(f)=f|_{A'}.
\]
Relația este cea indusă de $h$:
\[
 \sim\;=R_h,
 \qquad f\sim g\iff h(f)=h(g)
 \iff f|_{A'}=g|_{A'}.
\]
Prin urmare, aplicația indusă $\overline h$ este injectivă.

\medskip
\noindent Pentru orice aplicație $t:A'\longrightarrow B$ căutăm
$f:A\longrightarrow B$ astfel încât $f|_{A'}=t$.
Alegem $a\in B$ și definim
\[
 f(x)=
 \begin{cases}
   t(x),&x\in A',\\
   a,&x\in A\setminus A'.
 \end{cases}
\]
Aceasta este prelungirea indicată de diagramă:
\[
\begin{tikzcd}[column sep=large,row sep=large]
 A'\arrow[r,hook,"i"]\arrow[d,"t"']
    & A\arrow[dl,dashed,"f"]\\
 B &
\end{tikzcd}
\qquad f\circ i=t.
\]
Astfel, cum apare la sfârșitul paginii,
\[
 \forall\,g:A'\longrightarrow B,\quad
 \exists\,f:A\longrightarrow B
 \quad\text{a.î.}\quad f|_{A'}=g.
\]

\endgroup
; fără preambul sau pagină nouă.
% NOTE EDITORIALE:
% - Manuscrisul folosește Hom(A,B) pentru mulțimea tuturor aplicațiilor A -> B.
% - O clasă de echivalență este numită în notițe „germenele” lui f pe A'.
%   În terminologia uzuală, germenii presupun acordul pe o vecinătate;
%   relația de aici cere numai egalitatea restricțiilor la A'.
% - Formula exactă de după „germenele lui f” este neclară; pare să
%   trimită la A', posibil prin expresia „în punctul lui A'”.
% - Diagrama cu restricția este reprodusă folosind tikz-cd; p este
%   proiecția canonică, iar h aplicația de restricție.
% - Construcția de la final alege a în B, presupunând implicit B nevidă
%   dacă A\A' este nevidă; această ipoteză nu este enunțată în manuscris.
% - Semnele barate lângă diagrame nu sunt reconstituite.
% - Manuscrisul se oprește la enunțarea existenței unei prelungiri;
%   nu am adăugat o încheiere nouă a demonstrației.

\subsection*{Seminar 3 - Pagina 8}
\begingroup
\setlength{\abovedisplayskip}{5pt}
\setlength{\belowdisplayskip}{5pt}
\setlength{\abovedisplayshortskip}{3pt}
\setlength{\belowdisplayshortskip}{3pt}

Fie $A,B$ mulțimi și $A'\subseteq A$.
Pentru $f,g\in\operatorname{Hom}(A,B)$, definim
\[
 f\sim g\overset{\mathrm{def}}\iff f|_{A'}=g|_{A'}.
\]
Să se arate că $\sim$ este relație de echivalență și că
\[
 \operatorname{Hom}(A,B)/{\sim}
 \ \overset{\mathrm{bij}}{\cong}\ \operatorname{Hom}(A',B).
\]
Un element din mulțimea factor se numește în notițe
\emph{germenele} lui $f$ pe $A'$.

\medskip
\noindent\underline{Demonstrație.}
Considerăm proiecția canonică $p$ și aplicația de restricție $h$:
\[
\begin{tikzcd}[column sep=huge,row sep=large]
 \operatorname{Hom}(A,B)
    \arrow[r,"p"]\arrow[d,"h"']
  & \operatorname{Hom}(A,B)/{\sim}
       \arrow[dl,dashed,"\overline h"] \\
 \operatorname{Hom}(A',B) &
\end{tikzcd}
\qquad h(f)=f|_{A'}.
\]
Relația este cea indusă de $h$:
\[
 \sim\;=R_h,
 \qquad f\sim g\iff h(f)=h(g)
 \iff f|_{A'}=g|_{A'}.
\]
Prin urmare, aplicația indusă $\overline h$ este injectivă.

\medskip
\noindent Pentru orice aplicație $t:A'\longrightarrow B$ căutăm
$f:A\longrightarrow B$ astfel încât $f|_{A'}=t$.
Alegem $a\in B$ și definim
\[
 f(x)=
 \begin{cases}
   t(x),&x\in A',\\
   a,&x\in A\setminus A'.
 \end{cases}
\]
Aceasta este prelungirea indicată de diagramă:
\[
\begin{tikzcd}[column sep=large,row sep=large]
 A'\arrow[r,hook,"i"]\arrow[d,"t"']
    & A\arrow[dl,dashed,"f"]\\
 B &
\end{tikzcd}
\qquad f\circ i=t.
\]
Astfel, cum apare la sfârșitul paginii,
\[
 \forall\,g:A'\longrightarrow B,\quad
 \exists\,f:A\longrightarrow B
 \quad\text{a.î.}\quad f|_{A'}=g.
\]

\endgroup
; fără preambul sau pagină nouă.
% NOTE EDITORIALE:
% - Manuscrisul folosește Hom(A,B) pentru mulțimea tuturor aplicațiilor A -> B.
% - O clasă de echivalență este numită în notițe „germenele” lui f pe A'.
%   În terminologia uzuală, germenii presupun acordul pe o vecinătate;
%   relația de aici cere numai egalitatea restricțiilor la A'.
% - Formula exactă de după „germenele lui f” este neclară; pare să
%   trimită la A', posibil prin expresia „în punctul lui A'”.
% - Diagrama cu restricția este reprodusă folosind tikz-cd; p este
%   proiecția canonică, iar h aplicația de restricție.
% - Construcția de la final alege a în B, presupunând implicit B nevidă
%   dacă A\A' este nevidă; această ipoteză nu este enunțată în manuscris.
% - Semnele barate lângă diagrame nu sunt reconstituite.
% - Manuscrisul se oprește la enunțarea existenței unei prelungiri;
%   nu am adăugat o încheiere nouă a demonstrației.

\subsection*{Seminar 3 - Pagina 8}
\begingroup
\setlength{\abovedisplayskip}{5pt}
\setlength{\belowdisplayskip}{5pt}
\setlength{\abovedisplayshortskip}{3pt}
\setlength{\belowdisplayshortskip}{3pt}

Fie $A,B$ mulțimi și $A'\subseteq A$.
Pentru $f,g\in\operatorname{Hom}(A,B)$, definim
\[
 f\sim g\overset{\mathrm{def}}\iff f|_{A'}=g|_{A'}.
\]
Să se arate că $\sim$ este relație de echivalență și că
\[
 \operatorname{Hom}(A,B)/{\sim}
 \ \overset{\mathrm{bij}}{\cong}\ \operatorname{Hom}(A',B).
\]
Un element din mulțimea factor se numește în notițe
\emph{germenele} lui $f$ pe $A'$.

\medskip
\noindent\underline{Demonstrație.}
Considerăm proiecția canonică $p$ și aplicația de restricție $h$:
\[
\begin{tikzcd}[column sep=huge,row sep=large]
 \operatorname{Hom}(A,B)
    \arrow[r,"p"]\arrow[d,"h"']
  & \operatorname{Hom}(A,B)/{\sim}
       \arrow[dl,dashed,"\overline h"] \\
 \operatorname{Hom}(A',B) &
\end{tikzcd}
\qquad h(f)=f|_{A'}.
\]
Relația este cea indusă de $h$:
\[
 \sim\;=R_h,
 \qquad f\sim g\iff h(f)=h(g)
 \iff f|_{A'}=g|_{A'}.
\]
Prin urmare, aplicația indusă $\overline h$ este injectivă.

\medskip
\noindent Pentru orice aplicație $t:A'\longrightarrow B$ căutăm
$f:A\longrightarrow B$ astfel încât $f|_{A'}=t$.
Alegem $a\in B$ și definim
\[
 f(x)=
 \begin{cases}
   t(x),&x\in A',\\
   a,&x\in A\setminus A'.
 \end{cases}
\]
Aceasta este prelungirea indicată de diagramă:
\[
\begin{tikzcd}[column sep=large,row sep=large]
 A'\arrow[r,hook,"i"]\arrow[d,"t"']
    & A\arrow[dl,dashed,"f"]\\
 B &
\end{tikzcd}
\qquad f\circ i=t.
\]
Astfel, cum apare la sfârșitul paginii,
\[
 \forall\,g:A'\longrightarrow B,\quad
 \exists\,f:A\longrightarrow B
 \quad\text{a.î.}\quad f|_{A'}=g.
\]

\endgroup
